\subsection{Comparison of the recorded and mixed component volumes}
\label{mixrec:section}

Fix a marked cover with coefficient field \(F\), Galois CM host
field \(E\), level integer \(M'\), component \(Z\), and algebraic
weight \(u\). Set \(d_E=[E:\mathbf Q]\) and
\begin{align*}
 S_B&=\{p\text{ prime}:\ p\le B\ \text{or}\ p\mid 2|\Delta_E|M'\},\\
 c_p&=\frac{1}{p-1}+\frac{v_p(|\Delta_E|)}{d_E},
 \qquad K_B=\sum_{p\in S_B}c_p\log p .
\end{align*}
The recorded lattice imposes nonnegative valuations at all recorded
algebraic integral frames. The mixed lattice replaces these tests
over \(S_B\) by Hodge--Tate frames. Both use the same rational cusp
rules on \(Z\), and the same supremum over the good embeddings and
marked archimedean frames. Write \(V_{\rm rec}(t),V_{\rm mix}(t)\)
for their component volumes, counted on weights \(mu\) with
\(m\in8\mathbf Z_{>0}\). With the same \(\delta_Z=\operatorname{topDeg}_Z(u)>0\)
as in Proposition~\ref{zetageo:mean}, put
\(\ell_\star=\inf_t(V_\star(t)/\delta_Z-t)\).

Let \(a_s\) be the exponents of the tangent diagonal in
\(\code{exOf}(u,0)\), and let
\(W_+=\sum_s(a_s)_+\), \(W_-=\sum_s(-a_s)_+\).
At a CM-type embedding they are the first \(p_0\) entries of \(u\).
At its conjugate they are the negatives of the remaining \(q_0\)
entries in reverse order, where \((p_0,q_0)\) is the signature.
In particular \(W_++W_-=\|u\|_1\).

\begin{proposition}[Changing finite tests at a fixed weight]\label{mixrec:change}
For every \(\eta>0\) and \(t\in\mathbf R\),
\[
     V_{\rm rec}(t)\le V_{\rm mix}(t+W_+K_B+\eta),\qquad
     V_{\rm mix}(t)\le V_{\rm rec}(t+W_-K_B+\eta).
\]
The component slopes satisfy the extended-real inequalities
\[
           \ell_{\rm rec}-W_+K_B
                  \ \le\ \ell_{\rm mix}\
                  \le\ \ell_{\rm rec}+W_-K_B .
\]
These bounds hold also when a signature block is empty.
\end{proposition}
\begin{proof}
Fix \(p\in S_B\), a place \(\mathfrak l\) of \(F\) above it and a
Hodge--Tate test at \(\mathfrak l\). Its definition includes its
underlying algebraic integral frame \(y_0\), its ramification
normalization, and a lift to the \(p\)-adic comparison field. Thus
\(y_0\) defines a recorded test \(v_0\) at the same place. Conversely
every recorded test there has such a lift. Write \(w\) for the
Hodge--Tate test and \(e_{\mathfrak l}=\operatorname{ord}_{\mathfrak l}(p)\).
The lifted recorded frame and the Hodge--Tate frame differ by a
triangular matrix with tangent diagonal \(t_s\). In the multiplicative
valuation \(\nu\) of the comparison field one has
\[
       -c_p\le h_s:=\frac{\log\nu(t_s)}{\log\nu(p)}\le0 .
\]
Indeed the recorded eigenspace lattice is contained in the
Hodge--Tate lattice, giving \(\nu(t_s)\ge1\). The reverse inclusion
after multiplication by a \((p-1)\)-st root \(b\) of \(p\) and a
local different factor \(d\) gives \(\nu(bdt_s)\le1\). Put
\(L=-\log\nu(p)>0\). In \code{SumAbsExOfLe.lean}, lines 328--344,
the logarithmic bounds before the final relaxation are
\(\log\nu(b)=-L/(p-1)\) and
\(\log\nu(d)\ge-Lv_p(|\Delta_E|)/d_E\). Their sum gives
\(0\le\log\nu(t_s)\le Lc_p\). The comparison of weights
and evaluation in lines 346--364 computes the same diagonal factor
for every cusp rule \(f\) of weight \(mu\), giving
\[
             w(f)=v_0(f)+m e_{\mathfrak l}\sum_s a_s h_s .
\]
The scalar coordinate has valuation zero. If the evaluation vanishes,
both readings are \(+\infty\), and the equality still holds. The
finite displacement lies between \(-m e_{\mathfrak l}c_pW_+\) and
\(m e_{\mathfrak l}c_pW_-\). The reverse choice of a lift for a
prescribed recorded frame is \code{htCmp} in the same file,
lines 378--387. This gives both quantifiers over tests.

Put \(A_{\pm,m}=\prod_{p\in S_B}p^{\lceil m c_pW_\pm\rceil}\).
At each \(\mathfrak l\mid p\), a rational integer power \(p^j\)
adds exactly \(j e_{\mathfrak l}\) to the additive valuation.
Multiplication by \(A_{+,m}\) therefore sends the recorded lattice
injectively into the mixed lattice of the same weight. Multiplication
by \(A_{-,m}\) gives the reverse injection. Away from \(S_B\) all
the factors are units and the tests agree. At every archimedean frame
the corresponding norm is multiplied by \(A_{\pm,m}\), with
\[
        0\le \log A_{\pm,m}-mW_\pm K_B
                    \le \sum_{p\in S_B}\log p .
\]
Here \(B\), the cover and \(u\) are fixed before the degree tends to
infinity. For all sufficiently large multiples of eight, the upper
error is at most \(m\eta\). The two injections then compare the
counts in the indicated norm balls at that same degree.
Taking their normalized logarithmic upper limits gives the first
two inequalities. Dividing by \(\delta_Z\), shifting \(t\) in
the infimum and letting \(\eta\) decrease to zero gives the claims
for the slopes. The normalizing coefficient-field degree is the
same on both sides.
\end{proof}

On the selected hosts, \(M'\) has at most fifteen distinct prime
factors and \(s=\code{ell}(D)\ge R_E:=\log\operatorname{rd}(E)\ge0\).
Put \(Q=\lfloor B\rfloor\) at \(B=\code{bndB}(u,g)\). Here
\(B\ge4fr^2\ge288\), since \(f\ge m\ge2\) and
\(r=k+m-1\ge6\). The sum of \(\log p/(p-1)\) for \(p\le Q\)
is at most \(\log B+3\) by \code{SumLogDivPredLeReal.lean}, lines
12--14. Among the remaining primes of \(S_B\), those in
\code{Host.small} contribute at most \(4R_E/Q\) by
\code{BadCount.lean}, lines 18--33. There are at most seventeen
others by \code{BadPrimes.lean}, line 56. Each contributes at most
\(\log(Q+1)/Q\), since \(t\mapsto\log t/(t-1)\) decreases for
\(t>1\). The discriminant parts sum to \(R_E\), since \(S_B\)
contains every discriminant prime. Consequently
\begin{align*}
 K_B&\le\log B+3+17\,\frac{\log(Q+1)}{Q}
                 +\left(1+\frac4Q\right)R_E\\
    &\le 2(\log B+R_E).
\end{align*}
The second line uses \(Q\ge288\), so the term with coefficient
seventeen is less than one and \(\log B>4\). The first line also gives
\(K_B\le(1+O(1/\log B))(\log B+R_E)\) as \(B\) grows, uniformly
in \(R_E\ge0\). The first summand of \(\code{errX}(s,u,g)\)
is \((\log B+s)\|u\|_1\), and its remaining term is nonnegative. Thus
\[
                       W_\pm K_B\le2\,\code{errX}(s,u,g).
\]
The component average therefore raises the recorded constant \(C\)
to at most \(C+2\) in the mixed bound. The original estimate charges
\(40v_p(|\Delta_E|)/d_E\) locally and gives \(C+44\) for \(B\ge4\)
by \code{htCost\_le} in \code{BudgetOfSizeAndWeight.lean}, lines
71--108.

For a fixed admissible stage,
\code{heightAt\_of} in \code{HeightOfMeanBound.lean}, lines 85--95,
uses that mean only at \(u'=\code{uPrime}(u_1)\) and at the gap
supplied by its \code{BudgetAt} witness. Hence a common recorded
bound at those particular pairs already gives the mean used by that
application. In the actual tower the stage \(j\) has signature
\((k+j-1,m-j)\). The definition \code{Covers} in
\code{CovolXOfCovol.lean}, lines 220--223, takes \(1\le j<m\).
At the compact stage \(j=m\) the separate \code{LowDef} bound on
line 247 enters through the Gram estimate and its base comparison.

\begin{corollary}[Carrying recorded sections]\label{mixrec:carry}
Fix a component and a witness for \code{Mult} from \(u'\) to \(u_1\)
with cost \(c\). Write \(V^L_{u_1}\) for the Bergman volume on that
component. For every \(t\in\mathbf R\) and \(\eta>0\),
\[
       V_{{\rm rec},u'}(t)\le V^L_{u_1}(t+c+\eta).
\]
\end{corollary}
\begin{proof}
At degree \(8k\) multiply each recorded section \(f\) by
\(\alpha_k\theta^k\), where \(\theta\) is the carrying section of
\code{Mult}. The balanced scalar lemma applied to its prescribed
finite and archimedean data gives \(\alpha_k\ne0\) with
\(v_{\mathfrak l}(\alpha_k)\ge 8k e(\mathfrak l)\) at every
finite place. Recorded integrality of \(f\) now gives the finite
bound for the product at every recorded test. At the good embeddings
the carrying bound and the archimedean bound for \(\alpha_k\) give
the Bergman cutoff \(t+c+\eta\) once the scalar lemma's fixed
constant is at most \(8k\eta\). The generic value of \(\theta\)
is nonzero, so multiplication is injective on the component cusp
rules. Taking the normalized upper limit proves the assertion.
These are the finite and archimedean inequalities of
\code{WeightAbsEntrySumNonneg.lean}, lines 194--260, with the
recorded integrality bound inserted directly.
\end{proof}

In the component height argument of \code{HeightOfMeanBound.lean},
lines 48--81, a recorded lower slope \(s_Z\) at \(u'\) can therefore
be passed directly to the Bergman upper bound at \(u_1\). With the
same short generators, birationality, degree comparison and
cutoff \(c_s\), letting \(\eta\) decrease to zero gives
\[
         (1-\rho)(s_Z+c_s)_+\le h_Z+c_s+c .
\]
The carrying cost \(c\) is the only shift between the two volumes
in this application.

In Proposition~\ref{zetageo:mean} the parameter sections and
products can now be tested in the recorded lattices. If the required
rational cusp sections for the generic field basis exist, and their
fixed-degree recorded lattice spans the rational cusp space, a
single positive integer clears all their denominators. Its logarithm
contributes \(O(1/m)\) in the monomial count on this component.
If recorded cusp sections \(f_j\) of the parameter degree already
give a birational system, the generic basis may be chosen among
products of those sections and \(s_0\). Indeed, over the field of
the parameter ratios the \(f_j/s_0\) generate the finite function
field extension whenever the \(s_i\) lie in their span.
Monomials in these generators span the extension as a vector space.
Choose a basis of monomials and use \(s_0\) to put their numerators
in a common degree. The product hypotheses of
Proposition~\ref{zetageo:mean} then give their recorded integrality.
The remaining quantitative input is the mean of the recorded
parameter heights in the marked supremum norm, with the projective
degree and rational cusp-space identifications of that proposition.
The Petersson bounds in Corollary~\ref{cmeig:final} concern a fixed
good indefinite cover and its final weight. Applying them throughout
the chosen family still requires the supremum estimates and the
identifications on its modified flag components. The final row lattice
\(RR\) has its own gluing and metric in the slope comparison of
Section~\ref{geo:translate}.
